\documentclass[10pt,a4paper]{amsart}
\usepackage[margin=19mm]{geometry}
\usepackage{amsmath,amssymb,mathtools}
\usepackage[scr=rsfs,cal=euler]{mathalfa}
\usepackage{xfrac}
\usepackage[unicode,hidelinks,bookmarks=false]{hyperref}
\newcommand{\ie}{i.e.\ }
\newcommand{\cf}{cf.\ }
\newcommand{\resp}{respectively\ }
\newcommand{\Subsection}{\subsection}
\providecommand{\nobreakdash}{\nobreak\hbox{-}}
\setlength{\emergencystretch}{2em}
\multlinegap0pt
\usepackage{bm}
\usepackage{bbm}
\usepackage{dsfont}
\usepackage{xspace}
\usepackage{cancel}
\usepackage{mathtools}
\usepackage{amsthm}
\makeatletter
\@ifundefined{X}{\newtheorem{X}{X}}{}
\@ifundefined{lemma}{\newtheorem{lemma}[X]{Lemma}}{}
\@ifundefined{proposition}{\newtheorem{proposition}[X]{Proposition}}{}
\makeatother
\DeclareMathOperator{\disc}{Disc}
\DeclareMathOperator{\gl}{GL}
\DeclareMathOperator{\pp}{\mathbb{P}}
\newcommand{\syma}{\textnormal{Sym}^3(A^2)}
\title[Determining monogenity of pure cubic number fields using ell]{Determining monogenity of pure cubic number fields using elliptic curves}
\author{Jordi Gu\`ardia, Francesc Pedret}
\date{Source: arXiv:2505.06213v2 (CC BY 4.0)}
\begin{document}
\maketitle

\noindent\emph{Source and licence.} Determining monogenity of pure cubic number fields using elliptic curves. Version arXiv:2505.06213v2 (CC BY 4.0), distributed under the Creative Commons Attribution 4.0 International licence, as recorded in the arXiv licence metadata for this exact version and shown on the abstract page https://arxiv.org/abs/2505.06213v2. Full rights notice: licenses/arxiv-2505.06213-CC-BY-4.0.txt.\par\smallskip

\noindent\textit{Editorial scope.} Counted unit: the Lemma whose source label is \texttt{None} in the article (source0.tex lines 270--275, source label \texttt{None}), delivered together with the article's own complete original proof of it (source0.tex lines 276--300, one \texttt{proof} environment of 25 lines). No mathematics is written, repaired or re-derived: statement and proof are the article's own text line for line. Required context retained uncounted: proposition: GGS bijection (lines 245--252). Every definition, display and label of those lines is retained unchanged. Omitted from this package and not redistributed: 1--244, 253--269, 301--883. Nothing inside a retained excerpt is omitted or altered: every retained body is delivered exactly as the article's lines are written, so no omission receipt of any kind accompanies this package. Citations: the retained excerpts carry no citation command at all, so no bibliography entry of the article is transcribed and no reference is redistributed. Reference adaptation: none was needed and none was made; every reference inside the delivered unit resolves inside it, so no unresolved reference or citation is produced. Printed slips kept literal: no printed word, symbol, hypothesis or number of the counted statement or of its proof was altered. Layout and class: the article's own class is \texttt{amsart} with packages including graphicx, multirow, amsmath, amssymb, amsfonts, amsthm, mathrsfs, appendix, xcolor, textcomp, manyfoot, booktabs, algorithm, algorithmicx; the wrapper is the lane's pinned \texttt{amsart} editorial wrapper at 10pt on A4, which restates the theorem-like environments the retained text needs and adds the article's own macro definitions that the retained text needs (\texttt{\textbackslash disc}, \texttt{\textbackslash gl}, \texttt{\textbackslash pp}, \texttt{\textbackslash syma}), with extra wrapper packages (bm, bbm, dsfont, xspace, cancel, mathtools). The delivered numbering is the unit's own. Assets: no retained span contains a figure environment, a graphics-inclusion command, a PDF-page inclusion, a table or a TikZ picture; no image file is copied into this package, the package declares no asset dependency and no image file is redistributed with it. Licence: version arXiv:2505.06213v2 of this article is distributed under the Creative Commons Attribution 4.0 International licence, as recorded in the arXiv licence metadata for this exact version and shown on the abstract page https://arxiv.org/abs/2505.06213v2. A full rights notice is delivered at licenses/arxiv-2505.06213-CC-BY-4.0.txt. Characters: every retained excerpt is pure ASCII, so the delivered engine, XeLaTeX with its default Latin Modern text font, reports no missing character.
\begin{proposition}
    \label{proposition: GGS bijection}
    There is a discriminant-preserving bijection 
    \begin{alignat*}{4}
        \left\{\parbox{12em}{\textnormal{Isomorphism classes of rank 3            free } A \textnormal{-algebras}}\right\}&\longleftrightarrow&&\left\{\gl_2(A)-\textnormal{orbits on }\syma\right\}\\
        [B]&\longleftrightarrow&&\gl_2(A)\star I_{\mathcal{B}}(X,Y)
    \end{alignat*}
\end{proposition}

\begin{lemma} Let $C_{\mathcal{B}}$ the plane projective curve given by  $I_{\mathcal{B}}(X,Y) = Z^3$ on $\pp^2(K)$, where $\mathcal{B}$ is any $A$-basis of $B$.  Then
\begin{itemize}
    \item[a)]  $C_{\mathcal{B}}$ is smooth. 
\item[b)] For any other $A$-basis $\mathcal{B}'$ of $B$, the curves $C_{\mathcal{B}}$ and $C_{\mathcal{B}'}$ are isomorphic over $K$.
\end{itemize}
\end{lemma}

\begin{proof}
  \begin{itemize}
      \item[a)]
    Let $g(X,Y,Z) = I_{\mathcal{B}}(X,Y) - Z^3$ and suppose there is a singular point $P = (x:y:z)$ on $C_{\mathcal{B}}$. The partial derivatives of $g$ will vanish at $P$, so $z = 0$. 
    
    Suppose $y = 0$. Then we have the equality $I_{\mathcal{B}}(1,0) = 0$. By construction, this implies that $1,\omega_2,\omega_2^2$ are linearly dependent, so $K[\omega_2] \subsetneq L$. Since $3 = [L:K] = [L:K[\omega_2]][K[\omega_2]:K]$, we must have $[K[\omega_2] : K] = 1$, i.e. $\omega_2 \in K$. However, this implies that there exist $a, b \in A$ such that $a = b\omega_2$, contradicting the fact that $\{1,\omega_2,\omega_3\}$ is an $A$-basis of $B$ and, in particular, linearly independent. Thus, assume $y \neq 0$. In this case, $P$ is of the form $(x:1:0)$. The condition that $P$ is a singular point implies 
    $$
    I_{\mathcal{B}}(x,1) = \frac{\partial}{\partial x} I_{\mathcal{B}}(x,1) = 0,
    $$
    i.e. $x$ is a double root of $I(x,1)$. However, 
    $$
    \Delta(I_{\mathcal{B}}(x,1)) = \Delta(I_{\mathcal{B}}) \overset{\ref{proposition: GGS bijection}}{=} \disc(L/K) \neq 0,
    $$ 
    so $I_{\mathcal{B}}(X,1)$ cannot have a double root.
    \item[b)] If $\mathcal{B}'$ is another $A$-basis of $B$, then $I_{\mathcal{B}} = \gamma\star I_{\mathcal{B}'}$ for some $\gamma \in \gl_2(A)$. Thus, the change of variables given in affine coordinates by 
    $$
    (x, y)\mapsto (x,y)\cdot\gamma^{-1}:C_{\mathcal{B}'} \to C_{\mathcal{B}}
    $$ 
    is defined over $K$. Its inverse is the map 
    $$
    (x, y)\mapsto (x,y)\cdot\gamma:C_{\mathcal{B}} \to C_{\mathcal{B}'},
    $$ 
    which is also defined over $K$. \qedhere
\end{itemize}
\end{proof}

\end{document}
